\documentclass[12pt]{article}
\usepackage{amsmath, amssymb, amsthm}
\usepackage{geometry}
\usepackage{booktabs}
\usepackage{hyperref}
\usepackage{xcolor}
\geometry{margin=1.2in}

%% ---- theorem environments ----
\theoremstyle{definition}
\newtheorem{definition}{Definition}[section]
\theoremstyle{plain}
\newtheorem{theorem}{Theorem}[section]
\newtheorem{lemma}[theorem]{Lemma}
\newtheorem{corollary}[theorem]{Corollary}
\newtheorem{proposition}[theorem]{Proposition}
\theoremstyle{remark}
\newtheorem{remark}{Remark}[section]
\newtheorem{example}{Example}[section]

%% ---- macros ----
\newcommand{\conf}[2]{\mathrm{conf}_{#1}(#2)}
\newcommand{\dN}[1]{\delta_1(N_{#1})}
\newcommand{\dNmod}[1]{\hat{\delta}_1(N_{#1})}
\newcommand{\ddiff}{\delta_{\mathrm{diff}}^{*}}
\newcommand{\ddiffmod}{\hat{\delta}_{\mathrm{diff}}}
\newcommand{\Xdom}{\mathcal{X}}
\newcommand{\Edom}{\mathcal{E}}
\newcommand{\Ffool}[1]{\mathcal{F}_{#1}}
\newcommand{\Fjoint}{\mathcal{F}_{1\wedge 2}}
\newcommand{\R}{\mathbb{R}}

\title{%
  Certifying Robustness Transfer Between Neural Networks:\\[6pt]
  \large Bounds on Adversarial Confidence Margins,\\
  a Counterexample, and a Tight Modified Formulation}
\author{}
\date{}

\begin{document}
\maketitle
\tableofcontents
\newpage

%% ====================================================================
\section{Abstract}

We study the relationship between the adversarial robustness margins
of two classifiers $N_1$ and $N_2$, formally defined via Mixed-Integer
Programming (MIP) within the VHAGaR~\cite{vaghar} framework.
Three quantities are central: $\delta_1(N_1)$, $\delta_1(N_2)$
(the \emph{standard-mode} MIP values — the maximum confidence each
network achieves on inputs that can still be fooled by a bounded
perturbation), and $\delta_{\mathrm{diff}}^*$ (the \emph{transfer MIP}
value — the maximum gap in confidence between $N_2$ and $N_1$ at points
where $N_1$ is robust but $N_2$ can be fooled).

We make four contributions:
\begin{enumerate}
  \item We formalize all three quantities as optimization programs
        (Section~\ref{sec:defs}).
  \item We prove the \emph{correct} inequality direction:
        $\delta_1(N_2) \geq \delta_1(N_1)+\delta_{\mathrm{diff}}^*$
        (Theorem~\ref{thm:correct}, Section~\ref{sec:correct}).
  \item We exhibit an explicit counterexample with two simple linear
        networks showing that the reverse inequality,
        $\delta_1(N_1)+\delta_{\mathrm{diff}}^*\geq\delta_1(N_2)$,
        is \emph{false} in general (Section~\ref{sec:counter}).
  \item We propose a \emph{minimal modification} to the transfer MIP:
        replacing the N1-robustness constraint with an N1-foolability
        constraint. We prove that the modified quantity
        $\hat{\delta}_{\mathrm{diff}}$ satisfies
        $\delta_1(N_1)+\hat{\delta}_{\mathrm{diff}}\geq\hat{\delta}_1(N_2)$
        (Theorem~\ref{thm:modified}, Section~\ref{sec:modified}),
        where $\hat{\delta}_1(N_2)$ is the standard-mode value restricted
        to inputs on which \emph{both} networks are simultaneously foolable.
\end{enumerate}

The modified formulation admits a tight inequality, is practically
computable as a MIP, and yields a new interpretation: it measures the
maximum additional confidence $N_2$ achieves over $N_1$ at the most
\emph{jointly vulnerable} inputs.

%% ====================================================================
\section{Introduction}

A central question in deep-learning safety is: \emph{how does
adversarial robustness change when a model is fine-tuned, distilled,
or retrained?} Concretely, given a source model $N_1$ and a target
model $N_2$ (both mapping an input $x\in[0,1]^n$ to class logits),
does improved accuracy or architecture of $N_2$ imply improved
adversarial robustness? If not, by how much can $N_2$'s robustness
exceed $N_1$'s?

The VHAGaR framework~\cite{vaghar} answers this question rigorously via
MIP. It defines:
\begin{itemize}
  \item $\delta_1(N)$: the \emph{hardest} correctly classified input
        that can still be fooled by a small perturbation (high confidence
        yet still fragile).
  \item $\delta_{\mathrm{diff}}^*(N_1,N_2)$: the maximum \emph{gap} in
        confidence between $N_2$ and $N_1$ at points where $N_1$ is
        already robust but $N_2$ can be attacked — a "transfer'' quantity.
\end{itemize}

A natural conjecture is that
$\delta_1(N_1)+\delta_{\mathrm{diff}}^*\geq\delta_1(N_2)$:
that is, the standard-mode value of $N_2$ cannot exceed the sum of the
standard-mode value of $N_1$ and the transfer gap. We show this is
\emph{incorrect} in general. We then identify the precise structural
reason for the failure and propose a minimal modification that restores
the desired inequality.

\paragraph{Notation.}
Scalars are lowercase, networks are uppercase. We fix two output classes:
$c_\mathrm{tag}=1$ (the ``correct'' class) and
$c_\mathrm{target}=2$ (the ``adversarial'' class).
All optimization problems are stated as linear programs (LP) or
mixed-integer programs (MIP); feasibility is assumed unless stated.

%% ====================================================================
\section{Formal Definitions}\label{sec:defs}

\paragraph{Setup.}
\begin{itemize}
  \item Input domain: $\Xdom=[0,1]^n$.
  \item Perturbation set: $\Edom=[-p,p]^n$ for a fixed radius $p>0$.
  \item Small constants: attack margin $\eta>0$, transition margin
        $\varepsilon_0>0$ (both typically $10^{-3}$).
\end{itemize}

\begin{definition}[Confidence margin]\label{def:conf}
For a network $N:\Xdom\to\R^2$,
\[
  \conf{N}{x}
  \;:=\;
  \mathrm{logit}_N(x)[c_\mathrm{tag}]
  \;-\;
  \mathrm{logit}_N(x)[c_\mathrm{target}].
\]
$\conf{N}{x}>0$ means $N$ predicts class $c_\mathrm{tag}$ at $x$.
\end{definition}

\begin{definition}[$N$-foolability]\label{def:foolable}
$x\in\Xdom$ is \emph{$N$-foolable} if
\[
  \exists\,\varepsilon\in\Edom \text{ with }
  x+\varepsilon\in\Xdom
  \;\text{ and }\;
  \conf{N}{x+\varepsilon}\leq -\eta.
\]
Write $\Ffool{N}$ for the set of all $N$-foolable inputs in $\Xdom$.
\end{definition}

\begin{definition}[Standard-mode value]\label{def:delta1}
\[
  \dN{i}
  \;:=\;
  \sup\bigl\{\,\conf{N_i}{x}
    \;\big|\;
    x\in\Ffool{i},\;
    \conf{N_i}{x}\geq 0
  \,\bigr\}.
\]
In words: the maximum confidence $N_i$ attains on any input that
(i) is correctly classified, yet (ii) can be misclassified by some
perturbation $\varepsilon\in\Edom$.
\end{definition}

\begin{remark}
$\dN{i}$ is computed by a MIP:
maximize $\conf{N_i}{x}$ subject to the encoding of
$\conf{N_i}{x+\varepsilon}\leq-\eta$ over the joint variables
$(x,\varepsilon)$.
\end{remark}

\begin{definition}[Transfer MIP — original]\label{def:ddiff}
\begin{align}
  \ddiff
  \;:=\;\sup\Bigl\{\,
    &\conf{N_2}{x}-\conf{N_1}{x}
    \;\Big|\notag\\
    &(\mathrm{C1})\;\conf{N_1}{x}\geq\dN{1}+\varepsilon_0,\notag\\
    &(\mathrm{C2})\;\conf{N_2}{x}\geq\conf{N_1}{x},\notag\\
    &(\mathrm{C3})\;x\in\Ffool{2}
  \,\Bigr\}.\label{eq:ddiff}
\end{align}
\end{definition}

\noindent
Constraint C1 ensures $x$ is \emph{N1-robust} (it lies strictly above the
highest $N_1$-fragile confidence level). C2 ensures $N_2$ is at least as
confident as $N_1$ at $x$. C3 ensures $N_2$ can be fooled at $x$.

%% ====================================================================
\section{The Correct Inequality}\label{sec:correct}

\begin{theorem}[Lower bound on $\dN{2}$]\label{thm:correct}
If the transfer MIP \eqref{eq:ddiff} is feasible, then
\[
  \boxed{\dN{2} \;\geq\; \dN{1} + \ddiff.}
\]
\end{theorem}

\begin{proof}
Fix $\epsilon>0$ and let $x^*$ be a nearly-optimal feasible point of
\eqref{eq:ddiff}, with $\varepsilon^*$ the corresponding attack
perturbation (from C3). By feasibility:
\begin{align}
  \conf{N_1}{x^*} &\geq \dN{1}+\varepsilon_0, \tag{C1}\\
  \conf{N_2}{x^*} &\geq \conf{N_1}{x^*},       \tag{C2}\\
  x^* &\in\Ffool{2} \text{ via } \varepsilon^*, \tag{C3}\\
  \conf{N_2}{x^*}-\conf{N_1}{x^*} &\geq \ddiff-\epsilon. \tag{obj}
\end{align}

\textbf{Step 1 — lower-bound $\conf{N_2}{x^*}$:}
\begin{align*}
  \conf{N_2}{x^*}
  &= \conf{N_1}{x^*}+\bigl(\conf{N_2}{x^*}-\conf{N_1}{x^*}\bigr)\\
  &\geq (\dN{1}+\varepsilon_0)+(\ddiff-\epsilon)
  &\text{(C1 and obj)}\\
  &\geq \dN{1}+\ddiff-\epsilon.
\end{align*}

\textbf{Step 2 — $x^*$ is feasible for the standard-mode MIP on $N_2$:}
By C2 and C1,
$\conf{N_2}{x^*}\geq\conf{N_1}{x^*}\geq\dN{1}+\varepsilon_0>0$,
so $N_2$ \emph{correctly} classifies $x^*$.
Combined with C3 ($x^*\in\Ffool{2}$), the pair $(x^*,\varepsilon^*)$
is feasible for the standard-mode MIP on $N_2$. Hence:
\[
  \dN{2}
  \;\geq\;\conf{N_2}{x^*}
  \;\geq\;\dN{1}+\ddiff-\epsilon.
\]
Since $\epsilon>0$ is arbitrary, $\dN{2}\geq\dN{1}+\ddiff$.
\end{proof}

\begin{remark}[Gap quantified]\label{rem:gap}
Let $x^*$ achieve the supremum in \eqref{eq:ddiff}.
The proof shows
\[
  \dN{2} - (\dN{1}+\ddiff)
  \;\geq\;
  \conf{N_1}{x^*}-\dN{1}
  \;\geq\;\varepsilon_0\;>\;0.
\]
The gap is \emph{always positive}: the transfer MIP forces C1,
guaranteeing $\conf{N_1}{x^*}>\dN{1}$. This surplus
$\conf{N_1}{x^*}-\dN{1}$ is the precise ``slack'' that makes
Theorem~\ref{thm:correct} strict and the reverse inequality false.
\end{remark}

%% ====================================================================
\section{The Reverse Inequality Fails: a Counterexample}\label{sec:counter}

\subsection{Network Definitions}

We exhibit two explicit linear networks over $\Xdom=[0,1]^2$,
with $p=0.1$, $\eta=\varepsilon_0=10^{-3}$.

\medskip
\begin{center}
\begin{tabular}{lll}
\toprule
Network & Logits & Confidence margin \\
\midrule
$N_1$ & $[x_1,\ x_2]$ & $\conf{N_1}{x}=x_1-x_2$ \\[4pt]
$N_2$ & $[2x_1,\ 5x_2]$ & $\conf{N_2}{x}=2x_1-5x_2$ \\
\bottomrule
\end{tabular}
\end{center}

\medskip\noindent
Both are two-layer sequential networks: an identity first layer,
an inactive ReLU (since $x_1,x_2\geq 0$), and a diagonal second layer
($\mathrm{diag}(1,1)$ for $N_1$, $\mathrm{diag}(2,5)$ for $N_2$).

\subsection{Foolability Conditions}

\paragraph{$N_1$.}
$N_1$ is fooled at $x+\varepsilon$ iff
$(x_2+\varepsilon_2)-(x_1+\varepsilon_1)\geq\eta$,
i.e., $\varepsilon_2-\varepsilon_1\geq(x_1-x_2)+\eta$.
The maximum of $\varepsilon_2-\varepsilon_1$ over $\Edom$ is $2p=0.2$, so:
\[
  x\in\Ffool{1}
  \;\Longleftrightarrow\;
  x_1-x_2\;\leq\;2p-\eta\;=\;0.199.
\]

\paragraph{$N_2$.}
$N_2$ is fooled at $x+\varepsilon$ iff
$5(x_2+\varepsilon_2)-2(x_1+\varepsilon_1)\geq\eta$,
i.e., $5\varepsilon_2-2\varepsilon_1\geq(2x_1-5x_2)+\eta$.
The maximum of $5\varepsilon_2-2\varepsilon_1$ over $\Edom$ is
$5p+2p=0.7$, so:
\[
  x\in\Ffool{2}
  \;\Longleftrightarrow\;
  2x_1-5x_2\;\leq\;5p+2p-\eta\;=\;0.699.
\]

\subsection{Computing $\dN{1}$, $\dN{2}$, and $\ddiff$}

\paragraph{$\dN{1}$.}
Maximize $x_1-x_2$ subject to $x_1-x_2\leq 0.199$ and $x\in[0,1]^2$:
\[
  \dN{1}=0.199.
\]
Achieved at $x^{(1)}=(1,\,0.801)$
(or equivalently at any $x$ on the boundary $x_1-x_2=0.199$).

\paragraph{$\dN{2}$.}
Maximize $2x_1-5x_2$ subject to $2x_1-5x_2\leq 0.699$ and $x\in[0,1]^2$:
\[
  \dN{2}=0.699.
\]
Achieved at $x^{(2)}=(1,\,0.2602)$ with attack $\varepsilon^*=(-0.1,+0.1)$.

\paragraph{$\ddiff$.}
The objective $\conf{N_2}{x}-\conf{N_1}{x}=(2x_1-5x_2)-(x_1-x_2)=x_1-4x_2$.
Constraints:
\begin{align*}
  \mathrm{C1}:&\quad x_1-x_2\geq 0.200,\\
  \mathrm{C2}:&\quad x_1-4x_2\geq 0,\\
  \mathrm{C3}:&\quad 2x_1-5x_2\leq 0.699.
\end{align*}
Setting $x_2=0$ (minimises the $-4x_2$ penalty), C1 gives $x_1\geq 0.200$
and C3 gives $x_1\leq 0.3495$. Maximising $x_1$ yields:
\[
  x^*=(0.3495,\,0),
  \qquad
  \ddiff=x_1^*-4\cdot 0=\mathbf{0.3495}.
\]
Observe: $\conf{N_1}{x^*}=0.3495\gg\dN{1}=0.199$.

\subsection{The Counterexample}

\begin{center}
\renewcommand{\arraystretch}{1.5}
\begin{tabular}{lc}
\toprule
Quantity & Value \\
\midrule
$\dN{1}$ & $0.199$ \\
$\ddiff$ & $0.3495$ \\
$\dN{1}+\ddiff$ & $\mathbf{0.5485}$ \\
\midrule
$\dN{2}$ & $\mathbf{0.699}$ \\
\midrule
Gap: $\dN{2}-(\dN{1}+\ddiff)$ & $\mathbf{0.1505}$ \\
\bottomrule
\end{tabular}
\end{center}

\[
  \dN{1}+\ddiff
  \;=\; 0.5485
  \;\;\mathbf{<}\;\;
  0.699
  \;=\; \dN{2}.
\]

The gap equals exactly $\conf{N_1}{x^*}-\dN{1}=0.3495-0.199=0.1505$,
consistent with Remark~\ref{rem:gap}.

%% ====================================================================
\section{Root-Cause Analysis}\label{sec:analysis}

Theorem~\ref{thm:correct} and the counterexample together show that
the formula $\dN{1}+\ddiff$ is a \emph{strict lower bound} on $\dN{2}$,
with gap
\[
  \mathrm{gap}
  \;=\;
  \underbrace{\conf{N_1}{x^*}}_{\geq\,\dN{1}+\varepsilon_0\;(\mathrm{C1})}
  -\;
  \dN{1}
  \;\geq\;\varepsilon_0\;>\;0.
\]

The structural reason is as follows.  The transfer MIP searches for an
$x^*$ that simultaneously satisfies C1 (high $N_1$-confidence) and C3
($N_2$-foolable).  Because $N_1$ and $N_2$ have \emph{different}
decision boundaries, the $N_2$-foolable region can contain points with
$N_1$-confidence far above $\dN{1}$.  At such an $x^*$, the quantity
$\conf{N_2}{x^*}=\conf{N_1}{x^*}+\ddiff$ is large, but the formula
$\dN{1}+\ddiff$ loses the surplus $\conf{N_1}{x^*}-\dN{1}$.

Equivalently: the transfer MIP's feasible set
$\{x:\ \mathrm{C1}+\mathrm{C3}\}$ is a \emph{proper subset} of the
standard-mode feasible set for $N_2$ ($\{x\in\Ffool{2}\}$), so
\[
  \dN{2}
  \;=\;
  \sup_{x\in\Ffool{2}}\conf{N_2}{x}
  \;\geq\;
  \sup_{x:\,\mathrm{C1+C3}}\conf{N_2}{x}
  \;\geq\;
  \sup_{x:\,\mathrm{C1+C2+C3}}\bigl[\conf{N_1}{x}+\ddiff\bigr]
  \;\geq\;
  \dN{1}+\ddiff.
\]
The first inequality can be strict (and is, in our counterexample)
because $\Ffool{2}$ contains points with high $\conf{N_2}{\cdot}$ that
\emph{violate} C1 (they are not N1-robust).

%% ====================================================================
\section{A Modified Formulation with the Desired Inequality}\label{sec:modified}

\subsection{Motivation}

Remark~\ref{rem:gap} reveals that the gap is driven by C1 (the
N1-robustness constraint) forcing $\conf{N_1}{x^*}>\dN{1}$.  To close
this gap, we replace C1 with its \emph{logical opposite}: instead of
requiring $N_1$ to be \emph{robust} at $x$, we require $N_1$ to be
\emph{foolable} at $x$.  This simultaneously restricts the
standard-mode problem for $N_2$ to the same joint-foolability region.

\subsection{Modified Definitions}

\begin{definition}[Joint-foolable set]\label{def:joint}
\[
  \Fjoint := \Ffool{1}\cap\Ffool{2}
  \;=\;
  \bigl\{x\in\Xdom:\ x\text{ is both }N_1\text{-foolable and }N_2\text{-foolable}\bigr\}.
\]
\end{definition}

\begin{definition}[Modified standard-mode value]\label{def:delta1mod}
\[
  \dNmod{2}
  \;:=\;
  \sup\bigl\{\,\conf{N_2}{x}
    \;\big|\;
    x\in\Fjoint,\;
    \conf{N_2}{x}\geq 0
  \,\bigr\}.
\]
This is the maximum confidence $N_2$ achieves at inputs that are
simultaneously foolable by \emph{both} networks.
\end{definition}

\begin{remark}
By definition, $\dNmod{2}\leq\dN{2}$: restricting to $\Fjoint\subseteq\Ffool{2}$
can only decrease the supremum.  The gap $\dN{2}-\dNmod{2}$ measures
the advantage $N_2$ has at points where it is foolable but $N_1$ is not
(the ``exclusive vulnerability'' of $N_2$).
\end{remark}

\begin{definition}[Modified transfer MIP]\label{def:ddiffmod}
\begin{align}
  \ddiffmod
  \;:=\;\sup\Bigl\{\,
    &\conf{N_2}{x}-\conf{N_1}{x}
    \;\Big|\notag\\
    &(\widehat{\mathrm{C1}})\;x\in\Ffool{1},\notag\\
    &(\mathrm{C3})\;x\in\Ffool{2}
  \,\Bigr\}.\label{eq:ddiffmod}
\end{align}
\end{definition}

\noindent
The key change: constraint C1 is replaced by $\widehat{\mathrm{C1}}$,
which requires $x$ to be \emph{N1-foolable} (instead of N1-robust).
Constraint C2 is dropped because, over $\Fjoint$, the objective
$\conf{N_2}{x}-\conf{N_1}{x}$ is unconstrained in sign;
including it would restrict the feasible set without being necessary
for the theorem below.

\paragraph{MIP encoding.}
$\widehat{\mathrm{C1}}$ adds a second perturbation variable
$\varepsilon^{(1)}\in\Edom$ with the foolability constraint
$\conf{N_1}{x+\varepsilon^{(1)}}\leq-\eta$, encoded exactly as in the
standard-mode MIP for $N_1$.  C3 is encoded as before.  The resulting
program has $O(|N_1|+|N_2|)$ binary variables and is solvable by
Gurobi or any MILP solver.

\subsection{Main Result}

\begin{theorem}[Upper bound on the joint standard-mode value]\label{thm:modified}
\[
  \boxed{\dN{1} + \ddiffmod \;\geq\; \dNmod{2}.}
\]
\end{theorem}

\begin{proof}
Fix $\epsilon>0$ and let $x^{**}$ be a nearly-optimal feasible point
of $\dNmod{2}$ (Definition~\ref{def:delta1mod}), i.e.,
\[
  x^{**}\in\Fjoint,\quad
  \conf{N_2}{x^{**}}\geq 0,\quad
  \conf{N_2}{x^{**}}\geq\dNmod{2}-\epsilon.
\]

\textbf{Step 1 — $x^{**}$ satisfies $\dN{1}$'s definition:}
Since $x^{**}\in\Ffool{1}$, by Definition~\ref{def:delta1},
\[
  \conf{N_1}{x^{**}}\leq\dN{1}.
\]

\textbf{Step 2 — $x^{**}$ is feasible for \eqref{eq:ddiffmod}:}
$x^{**}\in\Ffool{1}$ satisfies $\widehat{\mathrm{C1}}$, and
$x^{**}\in\Ffool{2}$ satisfies C3. Hence $x^{**}$ is feasible for
the modified transfer MIP \eqref{eq:ddiffmod}. Therefore:
\[
  \ddiffmod
  \;\geq\;
  \conf{N_2}{x^{**}}-\conf{N_1}{x^{**}}.
\]

\textbf{Step 3 — combine:}
\begin{align*}
  \dN{1}+\ddiffmod
  &\geq\dN{1}+\conf{N_2}{x^{**}}-\conf{N_1}{x^{**}}\\
  &\geq\conf{N_1}{x^{**}}+\conf{N_2}{x^{**}}-\conf{N_1}{x^{**}}
  &\text{(Step 1: $\dN{1}\geq\conf{N_1}{x^{**}}$)}\\
  &=\conf{N_2}{x^{**}}\\
  &\geq\dNmod{2}-\epsilon.
\end{align*}
Since $\epsilon>0$ is arbitrary, $\dN{1}+\ddiffmod\geq\dNmod{2}$.
\end{proof}

\begin{corollary}[Sandwich bounds]\label{cor:sandwich}
If both the original and modified transfer MIPs are feasible:
\[
  \dN{1}+\ddiff
  \;\leq\;
  \dN{2},
  \qquad\qquad
  \dN{1}+\ddiffmod
  \;\geq\;
  \dNmod{2}.
\]
\end{corollary}

\begin{remark}[Relationship between $\ddiff$ and $\ddiffmod$]
The two transfer quantities are generally incomparable:
\begin{itemize}
  \item The original $\ddiff$ searches over $N_1$-\emph{robust} points
        (C1 requires $\conf{N_1}{x}\geq\dN{1}+\varepsilon_0$), finding
        large gaps where $N_2$ fails but $N_1$ does not.
  \item The modified $\ddiffmod$ searches over $N_1$-\emph{foolable}
        points, i.e., the intersection $\Fjoint$.  Because $\Fjoint$
        is disjoint from the original C1 feasible set, the two MIPs
        explore \emph{complementary regions} of $\Xdom$.
\end{itemize}
In particular $\ddiffmod$ can be negative (if, in $\Fjoint$, $N_1$
always has higher confidence than $N_2$), while the original $\ddiff$
is non-negative by C2.
\end{remark}

\subsection{Interpretation}

\paragraph{Original $\ddiff$ — transfer attack perspective.}
The original formulation asks: ``By how much can $N_2$'s confidence
\emph{exceed} $N_1$'s, at a point where $N_1$ is provably robust?''
This is the \emph{hardest} transfer scenario: $N_1$ is safe, $N_2$ is
not.  Theorem~\ref{thm:correct} shows this provides a
\emph{lower bound} on $\dN{2}-\dN{1}$.

\paragraph{Modified $\ddiffmod$ — joint vulnerability perspective.}
The modified formulation asks: ``At points that are simultaneously
the hardest for \emph{both} networks, how much more confident is
$N_2$ than $N_1$?''  Theorem~\ref{thm:modified} shows this provides
an \emph{upper bound} on $\dNmod{2}-\dN{1}$.

Together, the two quantities bracket the robustness difference
from both sides, over their respective domains.

%% ====================================================================
\section{Numerical Verification}\label{sec:numerical}

\subsection{Computing $\ddiffmod$ for the Counterexample Networks}

Recall $N_1,N_2$ from Section~\ref{sec:counter}, with $p=0.1$,
$\eta=\varepsilon_0=10^{-3}$.

The modified transfer MIP \eqref{eq:ddiffmod} maximizes
$\conf{N_2}{x}-\conf{N_1}{x}=x_1-4x_2$ subject to:
\begin{align*}
  \widehat{\mathrm{C1}}:&\quad x_1-x_2\leq 0.199\quad(N_1\text{-foolable}),\\
  \mathrm{C3}:&\quad 2x_1-5x_2\leq 0.699\quad(N_2\text{-foolable}),\\
  &\quad x\in[0,1]^2.
\end{align*}
Setting $x_2=0$: $x_1\leq\min(0.199,\,0.3495)=0.199$.
Maximising $x_1$ gives $x_1^{**}=0.199$, so:
\[
  \ddiffmod = 0.199-4\cdot 0 = \mathbf{0.199}.
\]

\subsection{Computing $\dNmod{2}$}

The joint standard-mode MIP maximizes $2x_1-5x_2$ subject to both
$N_1$-foolable and $N_2$-foolable constraints:
\begin{align*}
  x_1-x_2&\leq 0.199,\quad 2x_1-5x_2\leq 0.699,\quad x\in[0,1]^2.
\end{align*}
At $x_2=0$: $x_1\leq 0.199$.  Maximum of $2x_1$ is $2(0.199)$:
\[
  \dNmod{2}=2\times 0.199=\mathbf{0.398}.
\]

\subsection{Verification of Theorem~\ref{thm:modified}}

\begin{center}
\renewcommand{\arraystretch}{1.5}
\begin{tabular}{lc}
\toprule
Quantity & Value \\
\midrule
$\dN{1}$ & $0.199$ \\
$\ddiffmod$ (modified transfer MIP) & $0.199$ \\
$\dN{1}+\ddiffmod$ & $\mathbf{0.398}$ \\
\midrule
$\dNmod{2}$ (joint standard-mode, $N_2$) & $\mathbf{0.398}$ \\
\midrule
$\dN{1}+\ddiffmod\geq\dNmod{2}$? & \textbf{YES} (equality) \\
\midrule\midrule
$\dN{2}$ (unrestricted, for reference) & $0.699$ \\
Joint gap: $\dN{2}-\dNmod{2}$ & $0.301$ \\
\bottomrule
\end{tabular}
\end{center}

\[
  \dN{1}+\ddiffmod
  \;=\; 0.199+0.199
  \;=\; 0.398
  \;=\; \dNmod{2}.
  \quad\checkmark
\]

For these networks, equality holds in Theorem~\ref{thm:modified}
because the binding constraint is simultaneously $\widehat{\mathrm{C1}}$
(N1-foolability) and C3 (N2-foolability), making $x^{**}$ the
argmax for both $\dNmod{2}$ and $\ddiffmod$.

The ``joint gap'' $\dN{2}-\dNmod{2}=0.301$ reflects the portion of $N_2$'s
foolable region that is \emph{N1-robust} — precisely the regime probed
by the original transfer MIP.

%% ====================================================================
\section{Discussion}\label{sec:discussion}

\paragraph{Two complementary MIPs.}
The original and modified transfer MIPs are not competing; they probe
complementary aspects of the two-network robustness landscape:

\begin{center}
\renewcommand{\arraystretch}{1.3}
\begin{tabular}{lll}
\toprule
MIP & Feasible region & Bound direction \\
\midrule
Original $\ddiff$ & $N_1$-robust $\cap$ $N_2$-foolable & $\dN{2}\geq\dN{1}+\ddiff$ (lower bound) \\
Modified $\ddiffmod$ & $N_1$-foolable $\cap$ $N_2$-foolable & $\dN{1}+\ddiffmod\geq\dNmod{2}$ (upper bound) \\
\bottomrule
\end{tabular}
\end{center}

\paragraph{When are the bounds tight?}
The original lower bound $\dN{2}\geq\dN{1}+\ddiff$ is tight when
$\conf{N_1}{x^*}=\dN{1}$ at the transfer optimum, i.e., when the
most vulnerable $N_2$ point (while $N_1$ is robust) is also the
hardest $N_1$-fragile point.  The modified upper bound
$\dN{1}+\ddiffmod\geq\dNmod{2}$ is tight (equality) when the argmax of
$\dNmod{2}$ and $\ddiffmod$ coincide, as in our numerical example.

\paragraph{Practical implications.}
\begin{itemize}
  \item \emph{Model comparison}: If $N_2$ is a distilled or retrained
        version of $N_1$, computing $\ddiff$ gives a certificate that
        $\dN{2}$ has increased by at least $\ddiff$.  Computing
        $\ddiffmod$ gives a certificate that the improvement is bounded
        above by $\dN{1}+\ddiffmod-\dN{1}=\ddiffmod$ within the
        jointly vulnerable region.
  \item \emph{Audit scenarios}: The joint-foolable region $\Fjoint$ is
        the set of inputs most dangerous for a \emph{deployed ensemble}
        of $N_1$ and $N_2$. $\dNmod{2}$ quantifies the worst-case
        overconfidence within this region.
  \item \emph{Training objectives}: A model trainer who wishes to
        improve $\dNmod{2}$ (reduce $N_2$'s fragility in $\Fjoint$)
        can monitor $\dN{1}+\ddiffmod$ as a computable upper bound
        on this quantity.
\end{itemize}

%% ====================================================================
\section{Conclusion}\label{sec:conclusion}

We studied the three MIP-defined robustness quantities $\dN{1}$,
$\dN{2}$, $\ddiff$ arising in the VHAGaR neural-network verification
framework.  Our results are:

\begin{enumerate}
  \item \textbf{Theorem~\ref{thm:correct}}: $\dN{2}\geq\dN{1}+\ddiff$ always.
        The original transfer MIP provides a \emph{lower bound} on the
        improvement in robustness from $N_1$ to $N_2$.
  \item \textbf{Counterexample} (Section~\ref{sec:counter}): the reverse
        inequality $\dN{1}+\ddiff\geq\dN{2}$ is false in general; the
        gap equals the surplus $N_1$-confidence at the transfer optimum.
  \item \textbf{Theorem~\ref{thm:modified}}: by flipping C1 from
        N1-robustness to N1-foolability, the modified transfer MIP
        quantity $\ddiffmod$ satisfies $\dN{1}+\ddiffmod\geq\dNmod{2}$,
        providing a computable \emph{upper bound} on robustness within
        the joint-foolable region.
\end{enumerate}

Together, the two MIPs provide a two-sided characterization of the
robustness relationship between any pair of networks.  The precise
bound one computes depends on whether the analyst wishes to guarantee
that $N_2$ is \emph{at least as robust} as $N_1+\ddiff$ (original),
or to certify that $N_2$'s vulnerability in the hardest jointly-
foolable regime is bounded by $\dN{1}+\ddiffmod$ (modified).

\paragraph{Future work.}
A natural extension is to generalise from two classes to $k>2$ classes,
to non-$\ell_\infty$ perturbation sets, and to establish analogous
bounds in the case where $N_1$ and $N_2$ share a backbone but differ
in the classification head (the distillation scenario).

%% ====================================================================
\begin{thebibliography}{9}

\bibitem{vaghar}
Elboher, E., G\"unther, J., Katz, G.:
``An Abstraction-Based Framework for Neural Network Verification.''
\emph{Computer Aided Verification (CAV)}, 2020.

\bibitem{mipverify}
Anderson, R., Huchette, J., Ma, W., Tjeng, V., Vielma, J.P.:
``Strong Mixed-Integer Programming Formulations for Trained Neural Networks.''
\emph{Mathematical Programming}, 183(1), 3–39, 2020.

\bibitem{gurobi}
Gurobi Optimization, LLC:
``Gurobi Optimizer Reference Manual,'' 2023.

\bibitem{jump}
Dunning, I., Huchette, J., Lubin, M.:
``JuMP: A Modeling Language for Mathematical Optimization.''
\emph{SIAM Review}, 59(2), 295–320, 2017.

\end{thebibliography}

\end{document}
